\documentclass[11pt]{article}

\usepackage[T1]{fontenc}
\usepackage{lmodern}
\usepackage[margin=1in]{geometry}
\usepackage{amsmath,amssymb,amsthm,mathtools}
\usepackage{microtype}

\newtheorem{theorem}{Theorem}
\newtheorem{lemma}[theorem]{Lemma}
\newtheorem{corollary}[theorem]{Corollary}
\begin{document}

Fix any $m \ge 2$ and define
\[
 \lambda^{m}_{\alpha} \coloneq \max_{\left\{ I \mid I \text{ has } m \text{ machines} \right\}}\lambda_{\alpha}(I).
\]
\newline
$\lambda^{m}_{\alpha}$ is well defined, since: 
\begin{itemize}
     \item $m$ is fixed and thus bounded
     \item $n \le m^2$ for the relevant set of instances we maximize over. 
     \item We may assume, without loss of generality, that the job lengths are within a compact set $p_j \in [0,1]$
\end{itemize}
We are thus determining the maximum of a continous function over a compact set. 
Note, that we allow jobs of size \(p_j=0\) solely to work with a closed set.
If this is the case for some job $j$, we simply delete it from the instance. 

\begin{theorem}[Job Bound]\label{thm:job-bound}
Let $I$ be a $\lambda_{\alpha}^{m}$ maximizing instance.  Then
\[
  n\leq3m-2.
\]
\end{theorem}



\begin{proof}
Consider an arbitrary assigment of jobs from an instance $I$ onto its machines. We denote this as a partition of the jobs into $m$
sets $P(I)=(P_1,\ldots,P_m)$.
Write
\[
    L^{i}(P(I)):=\sum_{j\in P_i}p_j,
 \qquad L^{min}(P(I)):=\min_i L^i(P(I)),
 \qquad M_i(P(I)):=\max_{j\in P_i}p_j,
\]
where $M_i=0$ if $P_i$ is empty. 
We say that a schedule $S$ agrees with or corresponds to a partition $P(I)$ and vice versa, if the assigment of jobs to machines is the same for $P(I)$ and $\sigma(I)$.
We now put a condition on when this is the case for $P(I)$.
$P(I)$ corresponds to a list schedule, iff 

\begin{equation}
    L^{i}(P(I))-M_{i}(P(I)) \le L^{\min}(P(I)).
    \label{eq:list-characterization}
\end{equation}


Indeed, if a list produces $P(I)$, consider the last job assigned to
machine $i$.  Immediately before that assignment, machine $i$ has load
at least $L^{i}(P(I))-M_i(P(I))$ and is a least-loaded machine.  A machine of final
load $L^{min}(P(I))$ has current load at most $L^{min}(P(I))$, so necessarily $L^{i}(P(I))-M_{i}(P(I)) \le L^{min}(P(I))$ for
a list schedule.

Conversely, given a partition $P(I)$ such that \ref{eq:list-characterization} holds,
choose a largest job on every machine and place it last in that
machines processing order. Order every other job arbitrarily on its
machine. 
Every start time is at most
$L^i(P(I))-M_i(P(I))\leq L^{min}(P(I))$. Create a list by ordering the jobs in order
of nondecreasing start time. Break ties arbitrarily. We now schedule according to this list: 
\newline
Whenever a job of prescribed start time $s$ is selected from the list,
every machine with an unassigned job has current load at least $s$, and
every machine that already has all of its jobs assigned, has final load at least $ L^{min}(P(I))\geq s$.  Hence the
designated machine is currently least loaded, and the order is a
list agreeing with the partition $P(I)$.

Let $I$ be an instance with $m$ machines, that attains $\lambda_{\alpha}^m$

Choose a cost-maximizing list partition $P^{W}(I)$ and a minimizing
partition $P^*(I)$. A cost maximizing/minimzing list partition is a partition such that, 
there exists a cost maximizing/minimzing list schedule $S^*(I)$/$\sigma^{W}(I)$ that agrees with $P^*(I)$/$P^{W}(I)$. 
 

Form a bipartite undirected graph $G$: its left vertices are
the $m$ machines of $P^{*}(I)$, its right vertices are the $m$ machines of
$P^W(I)$, and every job $j$ corresponds to an edge, of weight $p_j$, incident to the two machine 
nodes to which it is assigned in $P^{W}$/$P^{*}$.

For a right vertex $v$, write
\[
 w_v:=L^v(P^W(I)),
 \qquad L^{min}(P^W(I)):=\min_v w_v,
 \qquad M_v(P^W(I)):=\max_{e\ni v} p_e.
\]
$M_v(P^W(I))$ is of course just the maximal processing time of a job assigned to machine $v$.
We call an edge $e_v$ \emph{tight}, if
\begin{enumerate}
     \item $ w_v-M_v(P^W(I))=L^{min}(P^{W}(I))$ \label{eq:critical}
     \item $e_v$ is the unique edge of weight $M_v$ incident with $v$
\end{enumerate}


Let $C$ be the set of tight edges.
We show that every cycle of $G$ contains an edge of $C$.  Suppose to the
contrary that a cycle $Z$ avoids $C$. Starting arbitrarily, give the successive edges of $Z$
alternating signs $d_e\in\{+1,-1\}$ and set $d_e=0$ for any edge not in the cycle.
We perturb the processing times:
\[
 p_e(t):=p_e+t d_e.
\]
For all sufficiently small $t$, every processing time remains
positive. We will denote the instance that results from these perturbations by
$I(t)$. At every vertex of the cycle, one incident edge increases
and one decreases.  Hence all loads in both partitions $P^{*}(I(t))$ and
$P^W(I(t))$ remain unchanged.

The partition $P^W(I(t))$ remains list-realizable for both sufficiently
small signs of $t$.  This follows from
\eqref{eq:list-characterization}.
If the inequality is tight and the largest incident edge is unique,
that edge is tight, hence it is not on $Z$ and remains
unchanged by perturbation. If the maximum $M_i(P^W(I))$ is attained by several incident edges, then
either an unchanged maximum remains (if it is not part of the cycle) and thus the maximum does not
fall below $M_v(P(I))$ after perturbation which implies that $w_v-M_v(P(I(t)))\leq L^{min}(P^{W}(I(t)))$ remains true, or the two cycle edges are all the
maximizers. In this case, one of them increases for either nonzero
sign of $t$.  In all cases
$w_v-M_v(P(I(t)))\leq L^{min}(P^{W}(I(t)))$ is true. 
Thus our pertubed instance and the partitions we chose for $I$ still correspond to list schedules.

Define
\[
Q(I)\coloneq \sum_{j=1}^{n} p^2_j
\]

The job-square sum  of the pertubed instance satisfies
\begin{equation}\label{eq:Q-perturbation}
 Q(I(t)) \coloneq \sum_{j=1}^{n} (p_j(t))^2 = \sum_{j=1}^{n} (p_j+t d_j)^2 = \sum_{j=1}^{n}p_j^2+ \sum_{i=1}^{n}2tp_jd_j + t^2 \sum_{i=1}^{n}d_j^2.
\end{equation}
For at least one of the two signs of every sufficiently small nonzero $t$, this
is strictly larger than $Q(I)$.

Let
\[
 A^W:=\sum_{i=1}^{m} L_i(P^{W}(I))^2,
 \qquad A^{*}:=\sum_{i=1}^{m} L_i(P^{*}(I))^2.
\]
Since we consider only instances that attain the nontrivial ratio
$\lambda^{m}_{\alpha}>1$, we must have
$A^{W}>A^{*}$.


Consider now a cost maximizing schedule $S^W(I)$ that agrees with the partition $P^W(I)$. 
aswell as  a optimal schedule $\sigma^{*}(I)$ that agrees with the partition $P^{*}(I)$
It holds for these, that

\[
\lambda_{\alpha}^{m}=
\frac{\gamma_{\alpha}(S^W(I))}{\gamma_{\alpha}(\sigma^{*}(I))}
<
\frac{\gamma_{\alpha}(S^W(I(t)))}{\gamma_{\alpha}(\sigma^{*}(I(t)))} \le \lambda_{\alpha}^m
\]

Where we choose in \eqref{eq:Q-perturbation} the sign that increases $Q(I(t))$.
The inequality follows from a expansion; see Appendix~\eqref{app:ratio-calculation}
This contradicts the definition of
$\lambda_{\alpha}^m$.  Therefore every cycle contains a tight
edge, or equivalently,
\begin{equation}\label{eq:forest}
 G-C\quad\text{is a forest}.
\end{equation}

There is at most one tight edge at each right vertex.  Moreover, a right
vertex of minimum load cannot have a positive tight edge.
\begin{equation}\label{eq:C-bound}
 |C|\leq m-1.
\end{equation}

The bipartite graph has $2m$ vertices.  By
\eqref{eq:forest}, its non-tight part has at most $2m-1$ edges.
Every job is one edge of $G$, so
\[
 n=|E(G-C)|+|C|
 \leq(2m-1)+(m-1)=3m-2.
\]
\end{proof}

For Appendix: 
\begin{equation} \label{app:ratio-calculation}
\begin{aligned}
&
\frac{A^W+(2\alpha-1)Q(t)}
     {A^*+(2\alpha-1)Q(t)}
-
\frac{A^W+(2\alpha-1)Q}
     {A^*+(2\alpha-1)Q}
\\[1mm]
&=
\frac{
\bigl(A^W+(2\alpha-1)Q(t)\bigr)
\bigl(A^*+(2\alpha-1)Q\bigr)}
{\bigl(A^*+(2\alpha-1)Q(t)\bigr)
 \bigl(A^*+(2\alpha-1)Q\bigr)}
-
\frac{
\bigl(A^W+(2\alpha-1)Q\bigr)
\bigl(A^*+(2\alpha-1)Q(t)\bigr)}
{\bigl(A^*+(2\alpha-1)Q(t)\bigr)
 \bigl(A^*+(2\alpha-1)Q\bigr)}
\\[1mm]
&=
\frac{
(2\alpha-1)A^WQ
+(2\alpha-1)A^*Q(t)
-(2\alpha-1)A^WQ(t)
-(2\alpha-1)A^*Q}
{\bigl(A^*+(2\alpha-1)Q(t)\bigr)
 \bigl(A^*+(2\alpha-1)Q\bigr)}
\\[1mm]
&=
\frac{
(2\alpha-1)
\bigl(Q(t)-Q\bigr)
\bigl(A^*-A^W\bigr)}
{\bigl(A^*+(2\alpha-1)Q(t)\bigr)
 \bigl(A^*+(2\alpha-1)Q\bigr)}
\\[1mm]
&=
\frac{
(1-2\alpha)
\bigl(Q(t)-Q\bigr)
\bigl(A^W-A^*\bigr)}
{\bigl(A^*+(2\alpha-1)Q(t)\bigr)
 \bigl(A^*+(2\alpha-1)Q\bigr)} > 0.
\end{aligned}
\end{equation}

The strict inequality holds because
\[
1-2\alpha>0,\qquad
Q(t)-Q>0,\qquad
A^W-A^*>0,
\]
and
\[
A^*+(2\alpha-1)Q(t)>0,
\qquad
A^W+(2\alpha-1)Q(t)>0.
\]
\end{document}
